% Imporditud: tools/import_eksperimendid.py
% Allikas: Eesti füüsikaolümpiaadi eksperimendi ülesannete
%   kogu (Rael Kalda, 2023), ülesanne Ü89, lahendus L89.
\setAuthor{EFO žürii}
\setRound{piirkonnavoor}
\setYear{2003}
\setNumber{G E1}
\setDifficulty{3}
\setTopic{dünaamika}
\setVahendid{raamat, mille kaaned pole väga libedad, ümmargune pliiats ja joonlaud}
\setPunktid{10}
\setAllikas{Eksperimendi kogu Ü89, lahendus lk 71}
\setLahendusseis{toores}

\prob{Libisemise hõõrdejõud}
Mitu korda erineb hõõrdejõud libisemisel hõõrdejõust veeremisel?

\hint

\solu
a) Idee: raamatust saab moodustada kaldpinna, mille kaldenurka saab muuta, tõstes või langetades raamatu üht serva. Asetame pliiatsi ükskord kaldpinnale, nii, et see saaks sealt alla veereda, teinekord nii, et pliiats saaks libiseda. Leiame minimaalse kaldenurgaga asendi (määrates raamatu tagumise otsa kõrguse h). Mõõtes ära kaldpinna otsa sellised minimaalsed kõrgused hv ja hl, mille puhul pliiats veel veereb (või libiseb), saame leida otsitava suhte (2 p). b) Arvutusvalemi tuletamine: hõõrdejõu valem F = mgh/L (1 p), kus h --- kaldpinna kõrgus ja L --- raamatu pikkus, arvutusvalem

Fv = hv , Fl hl

kus indeks v tähistab veeremist ja l --- libisemist (1 p). c) hv mõõtmine (1 p), mitme mõõtmise keskmise leidmine (1 p). d) hl mõõtmine (1 p), mitme mõõtmise keskmise leidmine (1 p). e) Otsitava suhte arvutamine (1 p). Lubatud on 50\% erinevus võrreldes korraldajate endi poolt antud katsevahendite komplekti juures saadud väärtusega. NB! Vältimaks vigu peaksid katse läbi viima üksteisest sõltumatult vähemalt kaks korraldajat. f) Veahinnang (1 p).

NB! Joonisel on objektid ebaproportsionaalselt lähedal.
\probend
